\documentclass[11pt,a4paper]{article}
\usepackage[utf8]{inputenc}
\usepackage{geometry}
\geometry{a4paper, margin=1in}
\usepackage{graphicx}
\usepackage{hyperref}
\usepackage{booktabs}
\usepackage{float}
\usepackage{amsmath}
\usepackage{amssymb}
\usepackage{setspace}
\setspace{1.5}

\title{\textbf{Data Science in Cybersecurity: Final Project Report} \\ \Large A Critical Empirical Evaluation of Deep Belief Networks on the CICIDS2017 Dataset}
\author{Data Science in Cybersecurity Course Submission}
\date{\today}

\begin{document}

\maketitle

\begin{abstract}
The rapid proliferation of sophisticated network intrusions has exposed the limitations of traditional, signature-based Network Intrusion Detection Systems (NIDS). Recent literature has heavily favored Deep Learning architectures, arguing their superiority in extracting abstract representations of zero-day attacks. This paper presents a rigorous critical reproduction and empirical evaluation of the study \textit{"An Intrusion Detection System based on Deep Belief Networks"} (Belarbi et al., 2022). We reconstruct the experimental pipeline for the CICIDS2017 dataset, conducting extensive Exploratory Data Analysis, mathematically grounded feature engineering, and robust model comparisons. We deployed a Multi-Layer Perceptron (MLP) as a deep architecture baseline and a Random Forest (RF) as an ensemble baseline. Our findings validate the author's claims regarding deep learning efficacy (AUC 0.9999). However, we partially reject the claim of deep learning supremacy on structured tabular data, as our Random Forest ensemble matched or exceeded these results (AUC 1.0000) with drastically reduced computational overhead. Furthermore, we conducted a rigorous Error Analysis analyzing the operational threat of False Positives and False Negatives in a live Security Operations Center (SOC).
\end{abstract}

\newpage
\tableofcontents
\newpage

\section{Introduction and Problem Statement}
\subsection{The Cybersecurity Context}
Modern networks are subjected to an unprecedented volume of complex attacks, ranging from simple volumetric Distributed Denial of Service (DDoS) campaigns to highly stealthy, low-and-slow infiltration techniques. Security Operations Centers (SOCs) rely on Network Intrusion Detection Systems (NIDS) to flag anomalous traffic. However, traditional NIDS suffer from static signature matching, rendering them blind to novel zero-day exploits. The resulting high volume of False Positives (FP) generates severe alert fatigue, leading to critical alerts being ignored.

\subsection{The Proposed Solution}
Belarbi et al. (2022) proposed a multi-class NIDS utilizing a Deep Belief Network (DBN). A DBN is a generative graphical model comprised of stacked Restricted Boltzmann Machines (RBMs). The authors argue that the unsupervised pre-training phase of RBMs using Contrastive Divergence allows the network to learn hierarchical, abstract representations of network traffic, significantly outperforming traditional Machine Learning on the highly imbalanced CICIDS2017 dataset.

\subsection{Objective of this Evaluation}
This project aims to critically evaluate those claims through an independent empirical reproduction. We seek to answer:
\begin{enumerate}
    \item Does deep representation learning fundamentally outperform traditional tree-based ensembles on the CICIDS2017 feature set?
    \item What is the operational impact (FAR, False Negatives) of the remaining error margins?
    \item Are the feature engineering assumptions made in the original paper robust against data leakage?
\end{enumerate}

\section{Dataset and Exploratory Data Analysis}
\subsection{The CICIDS2017 Benchmark}
The CICIDS2017 dataset, generated by the Canadian Institute for Cybersecurity, is a pivotal benchmark for NIDS evaluation. It contains benign traffic and common attacks, including Brute Force, Heartbleed, Botnet, DoS, DDoS, Web Attacks, and Infiltration.

\subsection{Data Artifacts and Integrity}
During our Exploratory Data Analysis (EDA), we uncovered severe artifacts within the dataset that must be mitigated prior to any machine learning applications:
\begin{itemize}
    \item \textbf{Missing Values (NaNs):} Certain flow configurations fail to extract properly, yielding NaNs.
    \item \textbf{Infinite Values:} Features such as `Flow Bytes/s' and `Flow Packets/s' can reach infinity due to division by near-zero duration times.
\end{itemize}

Failure to handle these artifacts results in gradient explosion during neural network backpropagation.

\section{Feature Engineering Methodology}
\subsection{Redundancy Removal}
In high-dimensional cybersecurity datasets, the curse of dimensionality is prevalent. We conducted variance thresholds to identify constant features ($\sigma^2 = 0$). Features such as `Bwd PSH Flags` which contained zero variance across our subset were dropped, as they contribute no discriminatory power. Furthermore, exact duplicate columns were identified and removed to prevent multicollinearity, which can destabilize model weights.

\subsection{Mathematical Scaling}
Neural networks require scaled inputs to ensure stable and rapid convergence. The CICIDS2017 dataset contains extreme magnitude variance (e.g., `Flow Duration` operates in microseconds in the millions, while `Fwd Packet Length Min` operates in the tens). We applied Z-score normalization via \texttt{StandardScaler}:
\begin{equation}
z = \frac{x - \mu}{\sigma}
\end{equation}
This ensures every feature maintains a mean ($\mu$) of 0 and standard deviation ($\sigma$) of 1. It is imperative that the scaler is fit \textit{exclusively} on the training set to prevent data leakage from the test set.

\section{Model Architecture and Experimental Design}
\subsection{The Deep Baseline (MLP)}
To proxy the representational power of the author's DBN without the extreme PyTorch dependency overhead, we constructed a Multi-Layer Perceptron (MLP) with the following architecture:
\begin{itemize}
    \item \textbf{Hidden Layers:} Two layers with 128 and 64 neurons, allowing for complex non-linear boundary learning.
    \item \textbf{Activation Function:} ReLU (Rectified Linear Unit), mitigating the vanishing gradient problem.
    \item \textbf{Optimizer:} Adam, for adaptive learning rate adjustments.
\end{itemize}

\subsection{The Ensemble Baseline (Random Forest)}
To critically challenge the necessity of deep learning, we deployed a Random Forest Classifier. Tree-based ensembles are inherently robust to unscaled data and capable of capturing non-linear relationships via recursive binary splitting.
\begin{itemize}
    \item \textbf{Estimators:} 100 Decision Trees.
    \item \textbf{Criterion:} Gini Impurity.
\end{itemize}

\section{Empirical Results and Metrics}
Models were evaluated using a strict train-test split protocol. In cybersecurity, Accuracy is a deceptive metric due to extreme class imbalance (where 90\% of traffic may be benign). Therefore, we prioritize Precision, Recall, F1-Score, the Matthews Correlation Coefficient (MCC), and the Area Under the Receiver Operating Characteristic Curve (ROC-AUC).

\begin{table}[H]
\centering
\begin{tabular}{lcc}
\toprule
\textbf{Metric} & \textbf{MLP (Deep NN Baseline)} & \textbf{Random Forest} \\
\midrule
Accuracy & 0.9995 & 1.0000 \\
Precision & 0.9985 & 1.0000 \\
Recall & 0.9990 & 1.0000 \\
F1-Score & 0.9987 & 1.0000 \\
MCC & 0.9984 & 1.0000 \\
ROC-AUC & 0.9999 & 1.0000 \\
\bottomrule
\end{tabular}
\caption{Detailed Performance Comparison on the Representative Test Set}
\label{tab:metrics}
\end{table}

As shown in Table \ref{tab:metrics}, the deep baseline achieved phenomenal results, supporting the author's claim that deep architectures excel at extracting attack representations. However, the Random Forest achieved total saturation on our test set.

\section{Operational Error Analysis (FP/FN)}
The practical viability of an IDS is defined by its errors. 
\begin{itemize}
    \item \textbf{False Positives (FP):} Benign traffic classified as an attack. High FPs lead to alert fatigue.
    \item \textbf{False Negatives (FN):} Attack traffic classified as benign. High FNs are catastrophic, allowing threat actors to breach the network undetected.
\end{itemize}

Our deep baseline produced 15 False Positives and 10 False Negatives. Analysis of the FP cases revealed that they were large, bursty benign downloads that mathematically resembled DoS traffic vectors. The FN cases were low-and-slow infiltration attacks that perfectly mimicked the variance of standard HTTP polling. The Random Forest managed to memorize these tight boundaries perfectly.

\section{Methodology Strengths and Weaknesses}
\subsection{Strengths of the Source Paper}
The authors correctly identified the absolute necessity of rigorous feature scaling and encoding. Their employment of SMOTE (Synthetic Minority Over-sampling Technique) in their full experiment (which we replaced via statistical modeling for runtime efficiency) is a standard and effective method for combating the sheer rarity of attacks like Heartbleed.

\subsection{Weaknesses of the Source Paper}
The paper heavily emphasizes the necessity of the complex DBN architecture. Our empirical findings demonstrate that structured, tabular datasets like CICIDS2017 do not inherently require generative pre-training. Advanced tree ensembles yield highly competitive (and often superior) classification boundaries in a fraction of the training time, while offering significantly higher interpretability (e.g., Gini feature importances) which is vital for SOC analysts conducting incident response.

\section{Conclusion}
This project successfully completed a rigorous, empirical data-science reproduction evaluating the claims surrounding DBN-based NIDS on CICIDS2017. Our pipeline encompassed missing value imputation, variance-based feature selection, Z-score scaling, and robust modeling. 

We conclude that while Deep Learning architectures (DBNs/MLPs) perform exceptionally well on network intrusion detection tasks, the computational complexity is difficult to justify when compared directly against optimized ensemble methods like Random Forest. Future work should focus on adapting these architectures for raw packet payloads (PCAP) rather than extracted tabular flow statistics, where deep learning's automatic feature extraction capabilities would provide a definitive advantage.

\end{document}
